\section{Computation}

The theory rests on a single local rule applied across the whole crystal, beat after beat. For that to
grow atoms, weather, and minds, the substrate must be able in principle to express any structure at all.
The guarantee that secures this is computational universality. A system is universal when it can carry out
any computation that can be carried out and simulate any other computer.

This section establishes that the
$\{5,3,4\}$ dodecagrid has this property, at the level of an external theorem and at the level of the
model's own running dynamics. It then draws the conclusion that the substrate is a computer rather than
merely like one, marks the limit that universal computation is goal-neutral, and closes with the
undecidability that comes packaged with the same power.

\subsection{Cellular automata}

A cellular automaton is a grid of identical cells, each holding one value from a small finite set, updating
together on a shared clock by one fixed rule that reads only a cell and its immediate neighbors. The
machine has four parts.
\begin{itemize}
  \item The cells are a regular array of sites.
  \item The states are values from a small finite set, where two states already suffice for everything of interest.
  \item The local rule is a single lookup from a cell and its neighbors to that cell's next state, the same rule everywhere.
  \item The beats are synchronous, with the whole grid stepping at once.
\end{itemize}
There is no program counter, no central memory, and no operator. A
quiescent state, the resting value such that a cell surrounded by quiescence stays quiescent, keeps a
starting pattern finite so that only the live region acts.

The lasting surprise is that such rules are not dull. On the line, an elementary rule with two states and
three neighbors has only $256$ possibilities. One of them, Rule 110, was proved universal by Cook. In the
plane, Conway's Game of Life turns a four-line neighbor-counting rule into gliders, guns, oscillators, and
working calculators.

The complexity is not in the rule. It is in the patterns the rule allows. A law
statable in one breath, applied locally an enormous number of times, produces moving objects, memory, and
full computation \cite{wolfram2020}. This is the permission slip for the theory. If a trivial rule on a
flat board yields gliders and calculators, a well-chosen rule on a vast curved crystal can yield particles,
forces, and minds. A complex world does not require a complex law.

The deepest lesson concerns what counts as a thing. A glider in the Game of Life is not a fixed set of
cells. It is a pattern that moves, re-forming one step over at each beat on different cells while keeping
its shape. In the theory every thing is this kind of object. Every particle and every self is a stable
traveling pattern the rule sustains, made of nothing but arrangement, while the cells beneath it come and
go. This is why a self can persist while none of its underlying cells do.

\subsection{The perception rule as a cellular automaton}

The model's dynamics is a cellular automaton on the nose, not by analogy. The cells are the vibes, the
sites of the $\{5,3,4\}$ dodecagrid, each with twelve touching neighbors. The states are the three tones
$+1$, $0$, and $-1$. The local rule is the perception rule. For a pair of neighboring tones it shares
opposite tones to peace, hops a charge into an empty neighbor, spawns a balanced pair from peace under the
arrow, and leaves same tones inert. The beat sweeps the edges once with a no-double-move guard, which is
exactly the synchronous tick. The quiescent state is peace, the tone $0$, doing nothing until touched. The
arrow is the only asymmetry, the creation bias that keeps the grid alive rather than relaxing to dead peace.

This is not a flat $256$-rule line and it is not the Game of Life. It is a three-state, twelve-neighbor
cellular automaton on a three-dimensional hyperbolic honeycomb. It is the same kind of object, and it
inherits the same lesson, that simple uniform local rules can produce open-ended complexity. The
neighborhood of a cell, the thing the rule reads, is defined through Margenstern's recursive splitting of
the hyperbolic tiling into a tree of nested regions addressed by Fibonacci arithmetic, which is what turns
an undrawable infinite plane into a navigable, computing object \cite{margenstern2008}.

Calling the
universe a cellular automaton does not mean it runs inside an outside machine. There is no operator and no
screen. The substrate is the automaton, and running the rule on the substrate is what existence does.

\subsection{Universality on hyperbolic tilings}

The strongest external result is Margenstern's. Cellular automata on hyperbolic tilings, the dodecagrid
$\{5,3,4\}$ included, are computationally universal \cite{margenstern2013}. This is a rigorous
construction, and the $\{5,3,4\}$ is the exact mesh of the theory, so the theorem lands directly on the
substrate.

Building a universal machine in hyperbolic space faces a special obstacle. There is no uniform zoom.
Hyperbolic space has no similarity transformations, so the usual trick of nesting a shrunken copy of a
machine is unavailable. Margenstern works around this with the railway model.

Picture a computer made of
trains. A single locomotive runs forever along the tracks. At junctions sit switches that can be flipped,
and the configuration of all switches encodes the state of a register machine. As the locomotive passes
through it reads switches and flips them, and that reading and flipping is the computation. Registers are
stored as track loops the train can lengthen or shorten. A train confined to a track cannot drift apart the
way free-flying signals do when hyperbolic lines diverge, so the railway keeps the computation literally on
the rails.

Lifted into hyperbolic three-space, the dodecahedral tiles each have twelve neighbors and can
route signals over bridges the plane lacks, which shrinks the machine to five states, against nine for the
two-dimensional pentagrid version. A five-state weakly universal automaton therefore lives on the exact
mesh. Weakly universal means it computes on an infinite ultimately periodic background rather than on a
blank quiet plane, which is the standard accepted sense for hyperbolic constructions.

\subsection{Universality on the live dynamics}

An external theorem is not enough. The theory needs universality on the model's own rule. The experiment
P44 supplies it. The local update is the ternary signed-majority rule, where a cell takes the sign of the
sum of its fill-weighted neighbor tones plus a bias. With both fills set to $-1$ and the bias set to $+1$,
that update is exactly NAND. Encode \texttt{true} as $+1$ and \texttt{false} as $-1$. The cell then
computes $\mathrm{out} = \operatorname{sign}(1 - a - b)$. The four input cases are as follows.

\begin{center}
\begin{tabular}{ccccc}
\hline
$a$ & $b$ & $1 - a - b$ & $\mathrm{out} = \operatorname{sign}(\cdot)$ & NAND \\
\hline
$+1$ & $+1$ & $-1$ & $-1$ (false) & T NAND T = F \\
$+1$ & $-1$ & $+1$ & $+1$ (true) & T NAND F = T \\
$-1$ & $+1$ & $+1$ & $+1$ (true) & F NAND T = T \\
$-1$ & $-1$ & $+3$ & $+1$ (true) & F NAND F = T \\
\hline
\end{tabular}
\end{center}

The output is false only when both inputs are true, which is exactly NAND. NAND is functionally complete,
so every Boolean function follows, and P44 builds NOT, AND, OR, XOR, and a full adder. The adder computes
$a + b + \mathrm{carry}$ correctly across all eight inputs, so real arithmetic runs on the rule. P44 then
reproduces the table of Rule 110 from rule-NANDs and evolves it non-trivially, tying the dynamics to a
known-universal system expressed entirely in its own gate.

The decisive part is stronger. The gates are not left as functions. Each is built as a real subgraph of
cells with symmetric ternary fills, with inputs and a positive bias clamped. The free cells run the model's
asynchronous signed-majority rule to a fixed point, and the output cells are read. NAND built this way
settles to the correct answer, and a six-gate XOR composed as $\mathrm{AND}(\mathrm{OR}(a,b),
\mathrm{NAND}(a,b))$ also settles correctly, since bus widths decrease downstream so each gate's margin
beats the feedback from the next layer.

Two worries are answered at once. The gates are run by the live
dynamics, not painted on as a separate circuit, and every output is a genuine fixed point stable under any
update order. A clamped-input gate is an attractor, not a decaying pulse, which defuses the concern that an
asynchronous, dissipative dynamics could not hold a stable logical value.

The experiment P141 closes the loop with a full Turing-universal machine running directly on the mesh. The
target is a two-counter Minsky machine, one of the smallest known universal models. The mapping uses only
the substrate's primitives. A register is a region of cells in the dodecagrid, and its value is the number
of charged cells in that region. Increment is creating a charge, one unit of the arrow. Decrement is
annihilating one charge. Test-zero is counting the charge.

P141 loads a multiplication program built from
increment and decrement-or-jump-if-zero, which is a universal instruction set, and computes
$R_2 = R_0 \times R_1$ with a temporary register preserving the multiplicand. The machine computes
$a \times b$ correctly across a spread of cases including a zero case, with registers and operations made
of nothing but tone and the arrow.

\subsection{The substrate is a computer}

Three independent angles converge.
\begin{itemize}
  \item Margenstern proves the $\{5,3,4\}$ class universal as a theorem.
  \item P44 shows functional completeness on the live dynamics, with NAND and a six-gate XOR settling as stable fixed points.
  \item P141 runs a full universal machine with charge as memory.
\end{itemize}
The conclusion is not that the universe
is like a computer. It is that the universe is one.

This converts a wild claim into a modest one. The wild claim is that a single rule on a crystal produces
atoms, galaxies, weather, and minds. The modest fact behind it is that the rule is universal, so any
pattern that can exist at all is a pattern the rule can host. If a working mind is a pattern a computer
could in principle run, the crystal can host it, because the crystal can run anything a computer can.

Universality is about capacity, not freedom of outcome. The rule can in principle express anything, while
the actual universe is fixed completely by the one rule, the single seed, and the deterministic growth that
follows. The capacity is unlimited and the outcome is one.

A clean boundary remains. Universality settles structure, never feeling. The theory treats experience as
the substance of the cells from the start, rather than as a product of computation, so the felt interior is
its own question and is not hidden inside the universality claim.

\subsection{Computation is goal-neutral}

Universality is necessary but not sufficient for mind. A universal computer can compute any function, but
it has no preference among them. It will compute a multiplication, a heat equation, or pure noise with
equal indifference. Universal computation is goal-neutral. It is raw capacity with no direction.

Intelligence needs direction. A problem is a gap between a desired state and a current state, and a
solution is a path that closes it. Solving a problem needs both the means, which is universal computation,
and the ends, which is a goal that picks one path over another. The substrate supplies the ends through the
arrow, the directional asymmetry the conserving rule selects toward.

So P44 and P141 establish the means
only. They show the substrate can compute anything. The bridge from computable to intentional is taken up
in the section on intelligence, which adds the arrow and the goal-directed search that aims this universal
computation at a target.

\subsection{Undecidability and the super-Turing question}

Power and limitation arrive together. The moment a system can simulate a universal machine, the halting
problem for that system becomes undecidable, since any Turing-complete substrate drags undecidability in
with it. This is not a flaw. It is the formal signature of a world where the only way to know the future is
to live it. The undecidability concerns the observer's predictive shortcut, not the system's evolution. The
rule still produces a perfectly definite future. There is simply no algorithm that fast-forwards to the
answer in general.

The cleanest undecidable problem living natively on the substrate is the domino problem. Given a finite set
of tile types with edge-matching rules, can copies cover the whole plane with no gaps and no overlaps,
every edge matching its neighbor. This is already undecidable in the Euclidean plane, proved by Berger in
1966. The hyperbolic version stayed open for nearly three decades until Margenstern proved, in early 2007,
that the domino problem is undecidable in the hyperbolic plane, with Kari proving the same result
independently by a different method within days \cite{margenstern2013}.

The proof has the classic shape.
Force the tiling to draw a computer, then ask whether the computer halts. Inside the heptagrid $\{7,3\}$
the matching rules force a rigid background, and on that scaffold a hierarchy of nested triangles forms a
grid that becomes the space-and-time diagram of a Turing machine. The tiles complete into a full tiling if
and only if the machine never halts. Since halting is undecidable, so is tiling.

Because the substrate is a
tiling, asking which local-matching patterns extend to cover the whole grid is the domino problem in
disguise, so undecidability is a native feature of the geometry the theory is built on rather than an
import from outside.

The frontier question is whether a hyperbolic substrate could ever compute past the Turing barrier.
Margenstern and Grigorieff sketch a model that does, on an infinigrid, a tiling by polygons with infinitely
many sides. One such cell has infinitely many neighbors, so in one update it can test a property across
infinitely many cells at once, evaluating a quantifier in a single tick. Chaining such cells decides every
level of the arithmetical hierarchy, which would be super-Turing computation.

The caveats are decisive and
kept clearly speculative here. The model uses infinite objects, a tile with infinitely many sides and a
cell reading infinitely many neighbors instantaneously, and such a thing cannot be built. It belongs to the
same idealized family as infinite-time Turing machines and black-hole computation, with rigorous internal
mathematics and no physical realizability \cite{thooft2014}. The dodecagrid has cells with twelve
neighbors, not infinitely many, so the infinigrid result does not apply to the substrate as built. The
right reading is a clean theorem on undecidability and a clearly labeled thought experiment on
hypercomputation, never the two confused.
